% ============================================================================
% Fichier  : annexes/Ann_D_textes.tex
% Version  : 1.0 -- novembre 2026
% ============================================================================

% Évite un titre de section isolé en bas de page.
\providecommand{\gardetitre}{\par\vskip 0pt plus 9\baselineskip\penalty-300\vskip 0pt plus -9\baselineskip}

\chapter{Textes et Normes}
\label{ann:textes}

\gardetitre
\section*{Degré de vérification des textes cités}

\begin{tableaucours}{L{4.6cm}L{3.6cm}L{3.8cm}}{%
    \tableauHeader{Texte} &
    \tableauHeader{Source utilisée} &
    \tableauHeader{Citation dans le cours}
}
Loi n°~2024/017 du 23~décembre~2024 & Copie certifiée conforme du
Secrétariat général de la Présidence & Articles cités mot pour mot \\
Loi n°~2010/012 du 21~décembre~2010 & Texte officiel & Articles 41 à 47
cités mot pour mot \\
Code de procédure pénale, articles 102 et 154 & D'après
Minka Mi Nguidjoï et al. (2026) & Paraphrasés, non cités \\
Convention de Malabo (2014) & Texte de l'Union africaine~; statut d'après
Minka Mi Nguidjoï et al. (2026) & Paraphrasée \\
Convention de Budapest (2001) & Conseil de l'Europe~; statut d'après
Minka Mi Nguidjoï et al. (2026) & Paraphrasée \\
RGPD (UE) 2016/679 & Journal officiel de l'Union européenne & Paraphrasé,
pour comparaison \\
Normes ISO et ISO/IEC & Sous licence & Paraphrasées, jamais
reproduites \\
Méthode PIA de la CNIL & Publication libre & Paraphrasée \\
\end{tableaucours}

\gardetitre
\section*{Textes internationaux et étrangers cités}

\begin{tableaucours}{L{4.4cm}L{2.2cm}L{5.4cm}}{%
    \tableauHeader{Texte} &
    \tableauHeader{Nature} &
    \tableauHeader{Place dans le cours}
}
Lignes directrices de l'OCDE (1980, révisées en 2013)~\cite{ocde1980} &
Recommandation & Origine des principes, Chapitre~\ref{chap:international} \\
Principes directeurs des Nations Unies (1990)~\cite{onu1990} & Résolution &
Chapitre~\ref{chap:international} \\
Convention~108 (1981) et protocole de 2018~\cite{coe1981conv108,coe2018conv108plus} &
Traité du Conseil de l'Europe & Chapitre~\ref{chap:international}~; fiches
pays \\
Convention de Malabo (2014)~\cite{ua2014malabo} & Traité de l'Union
africaine & Chapitre~\ref{chap:international} \\
Acte additionnel de la CEDEAO (2010)~\cite{cedeao2010acte} & Acte
communautaire & Comparaison, Chapitre~\ref{chap:international} \\
Directive n°~07/08-UEAC-133-CM-18 de la CEMAC (2008)~\cite{cemac2008} &
Directive communautaire & Chapitres~\ref{chap:loi2010} et~\ref{chap:secteurs} \\
Convention de Budapest (2001)~\cite{budapest2001} & Traité du Conseil de
l'Europe & Chapitres~\ref{chap:international}, \ref{chap:transferts}
et~\ref{chap:recherche} \\
RGPD (UE) 2016/679~\cite{ue2016rgpd} & Règlement européen &
Chapitre~\ref{chap:rgpd}~; double cote aux Chapitres~\ref{chap:transferts}
et~\ref{chap:violations} \\
Clauses contractuelles types de 2021~\cite{ue2021clauses} & Décision de la
Commission & Chapitre~\ref{chap:transferts} \\
Règlement (UE) 2024/1689 sur l'intelligence artificielle~\cite{ue2024ia} &
Règlement européen & Chapitre~\ref{chap:profilage} \\
Loi ivoirienne n°~2013-450~\cite{ci2013loi} & Loi étrangère & Fiche pays,
Chapitre~\ref{chap:international} \\
Loi marocaine n°~09-08~\cite{ma2009loi} & Loi étrangère & Fiche pays,
Chapitre~\ref{chap:international} \\
NIST Privacy Framework~1.0~\cite{nist2020pf} & Cadre volontaire &
Chapitre~\ref{chap:normes} \\
NIST AI RMF~1.0~\cite{nist2023airmf} & Cadre volontaire &
Chapitre~\ref{chap:profilage} \\
Stratégie continentale de l'Union africaine sur l'intelligence
artificielle (2024)~\cite{ua2024ia} & Orientation politique &
Chapitre~\ref{chap:profilage} \\
\end{tableaucours}

Aucun de ces textes n'est un critère obligatoire pour une organisation qui
ne traite que des données de personnes situées au Cameroun, à l'exception
des traités auxquels le Cameroun est Partie et des actes de la CEMAC. Le
Chapitre~\ref{chap:loi2010} donne la méthode pour choisir ses critères.

\gardetitre
\section*{Textes d'application attendus}

À la date du cours, l'auteur n'a pas trouvé de texte publié pour les
dispositions suivantes, qui renvoient à la voie réglementaire ou à un
référentiel de l'Autorité~: organisation de l'Autorité (art.~53)~;
modalités d'autorisation (art.~19, 32)~; critères d'évaluation du risque et
de l'évaluation d'impact (art.~25, 33)~; durées de conservation (art.~13,
28)~; délais d'examen des demandes (art.~47)~; référentiel des mesures
(art.~24, 27, 57)~; mécanisme de certification (art.~34)~; contrôle
(art.~35)~; textes spécifiques des secteurs (art.~4, 72). Le lecteur
vérifiera l'état de ces textes au moment où il lit.

\gardetitre
\section*{Correspondance des articles et des chapitres}

\begin{tableaucours}{L{4.0cm}L{8.2cm}}{%
    \tableauHeader{Articles} &
    \tableauHeader{Chapitres}
}
1 à 5 & \ref{chap:donnees} \\
6 à 18 & \ref{chap:principes} \\
19 & \ref{chap:loi} \\
20, 30, 31 & \ref{chap:soustraitants} \\
21, 37 à 47 & \ref{chap:principes}, \ref{chap:droits} \\
22 & \ref{chap:violations} \\
24 à 28 & \ref{chap:securite}, \ref{chap:programme} \\
29, 36 & \ref{chap:cartographie} \\
32 & \ref{chap:transferts} \\
33 & \ref{chap:risque} \\
34 & \ref{chap:loi}, \ref{chap:normes} \\
44, 65 & \ref{chap:profilage} \\
48 à 52, 72 & \ref{chap:donnees}, \ref{chap:secteurs} \\
53 à 71 & \ref{chap:loi} \\
73 & \ref{chap:prologue}, \ref{chap:recherche} \\
\end{tableaucours}
